\begin{center}
\LARGE
Minimal Involutive Bases

\end{center}

\begin{center}
\Large
Vladimir P. Gerdt

\end{center}

\noindent
Date:  July 17th (Wednesday)\\
Time:  11:30-12:00\\

\begin{center}
\large
Abstract
\end{center}

In this paper we study of the uniqueness properties of involutive polynomial bases which are
redundant Gr\"obner bases of the special form. The most general involutive algorithmic techniques
is based on the concept of involutive monomial division which allows one to separate all the
variables into multiplicative and non-multiplicative subsets. The separation gives thereby the
self-consistent computational procedure for constructing an involutive basis by performing
non-multiplicative prolongations and multiplicative reductions. Every specific involutive division
generates a particular form of involutive computational procedure. In addition to three involutive
divisions used by Thomas, Janet and Pommaret for analysis of partial differential equations we
introduce two new ones. These two divisions much as Thomas division do not depend on the order
of variables. We prove noetherity and continuity of the new divisions. Given noetherian and
continuous division, we present an algorithm for constructing of the minimal involutive basis for a
polynomial ideal. This minimal basis is uniquely defined for any admissible monomial ordering. 

